\section{USAJMO 2013}
        \begin{enumerate}
        	\item (\textit{USAJMO 2013}) Are there integers $a$ and $b$ such that $a^5b+3$ and $ab^5+3$ are both perfect cubes of integers?


 



\textbf{Solution}

No, such integers do not exist.  This shall be proven by contradiction, by showing that if $a^5b+3$ is a perfect cube then $ab^5+3$ cannot be.

 Remark that perfect cubes are always congruent to $0$, $1$, or $-1$ modulo $9$.  Therefore, if $a^5b+3\equiv 0,1,\text{ or} -1\pmod{9}$, then $a^5b\equiv 5,6,\text{ or }7\pmod{9}$.

 If $a^5b\equiv 6\pmod 9$, then note that $3|b$.  (This is because if $3|a$ then $a^5b\equiv 0\pmod 9$.)  Therefore $ab^5\equiv 0\pmod 9$ and $ab^5+3\equiv 3\pmod 9$, contradiction.

 Otherwise, either $a^5b\equiv 5\pmod 9$ or $a^5b\equiv 7\pmod 9$.  Note that since $a^6b^6$ is a perfect sixth power, and since neither $a$ nor $b$ contains a factor of $3$, $a^6b^6\equiv 1\pmod 9$.  If $a^5b\equiv 5\pmod 9$, then 
 \[a^6b^6\equiv (a^5b)(ab^5)\equiv 5ab^5\equiv 1\pmod 9\implies ab^5\equiv 2\pmod 9.\]  Similarly, if $a^5b\equiv 7\pmod 9$, then 
 \[a^6b^6\equiv (a^5b)(ab^5)\equiv 7ab^5\equiv 1\pmod 9\implies ab^5\equiv 4\pmod 9.\]  Therefore $ab^5+3\equiv 5,7\pmod 9$, contradiction.

 Therefore no such integers exist.

 Amkan2022

 



\textbf{Solution 2}

We shall prove that such integers do not exist via contradiction.Suppose that $a^5b + 3 = x^3$ and $ab^5 + 3 = y^3$ for integers x and y. Rearranging terms gives $a^5b = x^3 - 3$ and $ab^5 = y^3 - 3$. Solving for a and b (by first multiplying the equations together and taking the sixth root) gives a = $(x^3 - 3)^\frac{5}{24} (y^3 - 3)^\frac{-1}{24}$ and b = $(y^3 - 3)^\frac{5}{24} (x^3 - 3)^\frac{-1}{24}$. Consider a prime p in the prime factorization of $x^3 - 3$ and $y^3 - 3$. If it has power $r_1$ in $x^3 - 3$ and power $r_2$ in $y^3 - 3$, then $5r_1$ - $r_2$ is a multiple of 24 and $5r_2$ - $r_1$ also is a multiple of 24. 

 Adding and subtracting the divisions gives that $r_1$ - $r_2$ divides 12. (actually, $r_1 - r_2$ is a multiple of 4, as you can verify if $\{^{5r_1 - r_2 = 24}_{5r_2 - r_1 = 48}$. So the rest of the proof is invalid.) Because $5r_1$ - $r_2$ also divides 12, $4r_1$ divides 12 and thus $r_1$ divides 3. Repeating this trick for all primes in $x^3 - 3$, we see that $x^3 - 3$ is a perfect cube, say $q^3$. Then $x^3 - q^3 = 3,$ and $(x-q)(x^2 + xq + q^2) = 3$, so that $x - q = 1$ and $x^2 + xq + q^2 = 3$. Clearly, this system of equations has no integer solutions for $x$ or $q$, a contradiction, hence completing the proof.

 Therefore no such integers exist.

 



\textbf{Solution 3}

Let $a^5b+3=x^3$ and $ab^5+3=y^3$. Then, $a^5b=x^3-3$, $ab^5=y^3-3$, and 
 \[(ab)^6=(x^3-3)(y^3-3)\]Now take $\text{mod }9$ (recall that perfect cubes $\equiv -1,0,1\pmod{9}$ and perfect sixth powers $\equiv 0,1\pmod{9}$) on both sides. There are $3\times 3=9$ cases to consider on what values $\text{mod }9$ that $x^3$ and $y^3$ take. Checking these $9$ cases, we see that only $x^3\equiv y^3\equiv 0\pmod{9}$ or $x\equiv y\equiv 0\pmod{3}$ yield a valid residue $\text{mod }9$ (specifically, $(x^3-3)(y^3-3)\equiv 0\pmod{9}$). But this means that $3\mid ab$, so $729\mid (ab)^6$ so 
 \[729\mid (x^3-3)(y^3-3)\iff 729\mid (27x'^3-3)(27y'^3-3)\iff 81\mid (9x'^3-1)(9y'^3-1)\] contradiction.

 



\textbf{Solution 4}

If $a^5b+3$ is a perfect cube, then $a^5b$ can be one of $5,6,7 \pmod 9$, so $(a^5b)^5 = a^{25} b^5$ can be one of $5^5 \equiv 2$, $6^5 \equiv 0$, or $7^5 \equiv 4 \pmod 9$. If $a$ were divisible by $3$, we'd have $a^5 b \equiv 0 \pmod 9$, which we've ruled out. So $\gcd(a,9) = 1$, which means $a^6 \equiv 1 \pmod 9$, and therefore $a^{25} b^5 \equiv ab^5 \pmod 9$.

 We've shown that $a b^5$ can be one of $0, 2, 4 \pmod 9$, so $ab^5 + 3$ can be one of $3, 5, 7 \pmod 9$. None of these are possibilities for a perfect cube, so if $a^5b+3$ is a perfect cube, $ab^5+3$ cannot be.

 



\textbf{Solution 5}

As in previous solutions, notice $ab^5,a^5b \equiv 5,6,7 \pmod 9$. Now multiplying gives $a^6b^6$, which is only $0,1 \pmod 9$, so after testing all cases we find that $ab^5\equiv a^5b \equiv 6 \mod 9$. Then since $\phi (9) = 6$, $ab^5\equiv \frac{a}{b}\pmod 9$ and $a^5b \equiv \frac{b}{a}\pmod 9$ (Note that $a,b$ cannot be $0\pmod 9$). Thus we find that the inverse of $6$ is itself under modulo $9$, a contradiction.

 



\textbf{Solution 6}

I claim there are no such a or b such that both expressions are cubes.

 Assume to the contrary $a^5b +3$ and $ab^5 + 3$ are cubes. 

 \textbf{Lemma 1}: If $a^5b +3$ and $ab^5 + 3$ are cubes, then $ab^5, a^5b \equiv 5,7 \pmod 9$

 \textbf{Proof} Since cubes are congruent to any of $0, 1, -1 \pmod 9$, $ab^5,a^5b \equiv 5,6,7 \pmod 9$. But if $ab^5 \equiv 6 \pmod 9$, $3|a$, so $a^5b \equiv 0 \pmod 9$, contradiction. A similar argument can be made for $ab^5 \neq 6 \pmod 9$. 

 
\textbf{Lemma 2}: If k is a perfect 6th power, then $k \equiv 0,1 \pmod 9$ 

 \textbf{Proof}: Since cubes are congruent to $0, 1,  -1 \pmod 9$, we can square, and get 6th powers are congruent to $0, 1 \pmod 9$.

 
Since $ab^5 \cdot a^5b = a^6 b^6 = (ab)^6$, which is a perfect 6th power, by lemma 2, $ab^5 \cdot a^5b \equiv 0,1 \pmod 9$.

 But, by lemma 1, $ab^5 \cdot a^5b \equiv 5 \cdot 5, 5 \cdot 7, 7 \cdot 7 \equiv 7, 8, 4 \pmod 9$. 

 So, $ab^5 \cdot a^5b$, which is an integer, can't go into any of the possible residue classes modulo 9, without breaking one of these lemmas. This is a contradiction, and the proof is complete. $\blacksquare$

 -AlexLikeMath

 


 



\textbf{Solution 7}

Note that $x^3 \equiv 0,1,-1 \pmod 9$. So $a^5b+3, ab^5+3\equiv 0,1,-1\pmod{9}$$\Rightarrow a^5b, ab^5 \equiv 5,6,7 \pmod 9$. 

 \begin{verbatim}
		$a^5b \equiv 5 \pmod 9 \;\;\;\;\;\; ab^5 \equiv 5 \pmod 9$           $a^5b \equiv 6 \pmod 9 \;\;\;\;\;\; ab^5 \equiv 6 \pmod 9$           $a^5b \equiv 7 \pmod 9 \;\;\;\;\;\; ab^5 \equiv 7 \pmod 9$.
\end{verbatim}

Multiplying $a^5b$ and $ab^5$ we get $(ab)^6$. By substituting values for $a^5b$ and $ab^5$ we can see that $(ab)^6 \equiv 0,1,3,6,7,8 \pmod 9$ but because cubes are not $3,6,7 \pmod 9$, and squares are not $-1\pmod 9$ so $\boxed{(ab)^6 \equiv 0,1 \pmod 9}$.

 
                \begin{itemize}
                    \item Case I: $(ab)^6\equiv 0 \pmod 9$.This is only true when $3|ab$. If $3|a,b$ then $a^5b+3 \equiv 3 \pmod 9$, which is not a perfect cube. The same can be seen with if $3|a$ then $a^5b+3 \equiv 3 \pmod 9$ and if $3|a$ then $b^5a+3 \equiv 3 \pmod 9$. So 6 mod 9 case is eliminated from the above relations of congruences\item Case II: $(ab)^6\equiv 1 \pmod 9$ for this we'll take a look at the above information regarding $a^5b$ and $ab^5$. If $a^5b \equiv 5 \pmod 9$ then,
                \end{itemize}


                \begin{verbatim}
		$(ab)^6=a^5b \cdot ab^5 \equiv 5(ab^5) \pmod 9 \Rightarrow ab^5 \equiv 2 \pmod 9 \Rightarrow ab^5+3 \equiv               5 \pmod 9$ 
\end{verbatim}

which is not possible. Similarly, if $a^5b \equiv 7 \pmod 9$ then,

 \begin{verbatim}
		$(ab)^6=a^5b \cdot ab^5 \equiv 7(ab^5) \pmod 9 \Rightarrow ab^5 \equiv 4 \pmod 9 \Rightarrow ab^5+3 \equiv              7 \pmod 9$
\end{verbatim}

which is not possible. So, no such integers.

 -Lakshya Pamecha

 



\textbf{Solution 8}

We proceed by contradiction. Assuming $a^5b + 3 = x^3$ for some integer $x$ and $ab^5 + 3 = y^3$ for another integer $y$, we can multiply the two equations to get $a^6b^6 + 3a^5b + 3ab^5 + 9 = x^3y^3$. We can rearrange this to $3a^5b+3ab^5 = x^3y^3 - a^6b^6 +9$.

 Looking at this $\pmod 9$, this becomes $3a^5b+3ab^5 \equiv x^3y^3-a^6b^6 \pmod 9$. Looking at the right hand side, the first term is either -1, 0, or 1 $\pmod 9$. The second is  0 or 1 $\pmod 9$, because it is a square of a cube. Additionally, since the LHS is divisible by three, so is the RHS.

 The only achievable value for the LHS that fits this is 0. So, $3a^5b+3ab^5 \equiv 0 \pmod 9$. Dividing by three, making sure to divide the modulus as well, $a^5b + ab^5 \equiv 0 \pmod 3$. Noticing $a^5b \equiv ab \pmod 3$ and $ab^5 \equiv ab \pmod 3$, we get, $2ab \equiv 0 \pmod 3$.

 This means either $3|a$ or $3|b$, or both. WLOG, $3|a$. Looking at the original expression mod 9, we get $a^5b + 3 \equiv -1, 0$,or $1 \pmod 9$. Since $a$ is divisible by three, $a^5$ is divisible by 243, so it is divisible by 9. This reduces the equation to $3 \equiv -1, 0$,or $1 \pmod 9$, which obviously can't be true, a contradiction. Therefore, our original assumption, that both expressions were perfect cubes, is incorrect.

 So, no, no such integers exist.

 -sillybone

 \textbf{Note} After you get that $a^5 b \equiv 5, 6, \text{or} 7 \pmod{9}$, you can also bash by listing all possible residues for $a^5$ and their corresponding values for $b$. At that point, you can easily see that all the solutions for $a$ and $b$ do not work in the other equation, meaning that there are no solutions that satisfy both equations.

 


	\item (\textit{USAJMO 2013}) Each cell of an $m\times n$ board is filled with some nonnegative integer.  Two numbers in the filling are said to be \textit{adjacent} if their cells share a common side.  (Note that two numbers in cells that share only a corner are not adjacent).  The filling is called a \textit{garden} if it satisfies the following two conditions:


 (i) The difference between any two adjacent numbers is either $0$ or $1$.


 (ii) If a number is less than or equal to all of its adjacent numbers, then it is equal to $0$.


 Determine the number of distinct gardens in terms of $m$ and $n$.


 



\textbf{Solution}

We claim that any configuration of $0$'s produces a distinct garden.  To verify this claim, we show that, for any cell that is nonzero, the value of that cell is its distance away from the nearest zero, where distance means the shortest chain of adjacent cells connecting two cells.  Now, since we know that any cell with a nonzero value must have a cell adjacent to it that is less than its value, there is a path that goes from this cell to the $0$ that is decreasing, which means that the value of the cell must be its distance from the $0 \rightarrow$ as the path must end.  From this, we realize that, for any configuration of $0$'s, the value of each of the cells is simply its distance from the nearest $0$, and therefore one garden is produced for every configuration of $0$'s.  

 
However, we also note that there must be at least one $0$ in the garden, as otherwise the smallest number in the garden, which is less than or equal to all of its neighbors, is $>0$, which violates condition $(ii)$.  There are $2^{mn}$ possible configurations of $0$ and not $0$ in the garden, one of which has no $0$'s, so our total amount of configurations is $\boxed{2^{mn} -1}$

 \textbf{Similar Second Solution and Inspiration} Note: "bordering" and "surrounding" mean that two cells have to touch on a side; one vertex is not good enough.

 We note that there is no real step to begin the problem, so start by constructing the base case: put ALL the zeroes into the board. Then note that all squares bordering it (corners alone don't count) have to be $1$. After doing that, what can a cell bordering a 1 have as its value? Either 1 or 2, based on (i). But if it were 1, then all cells surrounding it would have at least 1 as their value by (i). And by (ii) the value would have to be 0, which contradicts our initial construction. Therefore, all the cells bordering a 1 have to be 2. Looks too simple for JMO, right?We apply the same logic to a cell bordering value $n$: either it is $n$ or $n+1$. If it is $n$, however, by constraint (i) and (ii) we realize that this cell has to be 0 (because neighbors cannot be less than it), contradicting our construction again! Therefore, that cell has to be $n+1$. And so on until the grid is filled. Basically the problem reduces to finding all 0s, surrounding them with 1s, and surrounding the 1s with 2s, etc. Beautiful! I have established a bijection between the zeros placed in the grid and the arrangement, therefore there are $2^{mn}$ solutions right? NO! Take the smallest cell. It has to be $0$ via criterion (ii). So a case with zero zeros is apocryphal. Our answer is $2^{mn}-1$.

 


	\item (\textit{USAJMO 2013}) In triangle $ABC$, points $P,Q,R$ lie on sides $BC,CA,AB$ respectively.  Let $\omega_A$, $\omega_B$, $\omega_C$ denote the circumcircles of triangles $AQR$, $BRP$, $CPQ$, respectively.  Given the fact that segment $AP$ intersects $\omega_A$, $\omega_B$, $\omega_C$ again at $X,Y,Z$ respectively, prove that $YX/XZ=BP/PC$.


 



\textbf{Solution 1}


                    \begin{figure}[H]
                        \centering
                        \includegraphics[width=0.5\textwidth]{Solution/USAJMO_I_2013_3_Solution_Figure_1.png}
                    \end{figure}
                    

 In this solution, all lengths and angles are directed.

 Firstly, it is easy to see by that $\omega_A, \omega_B, \omega_C$ concur at a point $M$ (the Miquel point). Let $XM$ meet $\omega_B, \omega_C$ again at $D$ and $E$, respectively. Then by Power of a Point, we have 
 \[XM \cdot XE = XZ \cdot XP \quad\text{and}\quad XM \cdot XD = XY \cdot XP\] Thusly 
 \[\frac{XY}{XZ} = \frac{XD}{XE}\] But we claim that $\triangle XDP \sim \triangle PBM$. Indeed, 
 \[\measuredangle XDP = \measuredangle MDP = \measuredangle MBP = - \measuredangle PBM\] and 
 \[\measuredangle DXP = \measuredangle MXY = \measuredangle MXA = \measuredangle MRA = 180^\circ - \measuredangle MRB = \measuredangle MPB = -\measuredangle BPM\]Therefore, $\frac{XD}{XP} = \frac{PB}{PM}$. Analogously we find that $\frac{XE}{XP} = \frac{PC}{PM}$ and we are done.

 
courtesy v\_enhance, minor clarification by integralarefun

 



\textbf{Solution 2}

\href{https://www.flickr.com/photos/127013945@N03/14800492500/lightbox/}{Diagram}Refer to the Diagram link.

 By Miquel's Theorem, there exists a point at which $\omega_A, \omega_B, \omega_C$ intersect. We denote this point by $M.$ Now, we angle chase:
 \[\angle YMX = 180^{\circ} - \angle YXM - \angle XYM\]
 \[= 180^{\circ} - \angle AXM - \angle PYM\]
 \[= \left(180^{\circ} - \angle ARM\right) - \angle PRM\]
 \[= \angle BRM - \angle PRM\]
 \[= \angle BRP = \angle BMP.\]In addition, we have
 \[\angle ZMX = 180^{\circ} - \angle MZY - \angle ZYM - \angle YMX\]
 \[= 180^{\circ} - \angle MZP - \angle PYM - \angle BMP\]
 \[= 180^{\circ} - \angle MCP - \angle PBM - \angle BMP\]
 \[= \left(180^{\circ} - \angle PBM - \angle BMP\right) - \angle MCP\]
 \[= \angle BPM - \angle MCP\]
 \[= 180^{\circ} - \angle MPC - \angle MCP\]
 \[= \angle CMP.\]Now, by the Ratio Lemma, we have
 \[\frac{XY}{XZ} = \frac{MY}{MZ} \cdot \frac{\sin \angle YMX}{\sin \angle ZMX}\]
 \[= \frac{\sin \angle YZM}{\sin \angle ZYM} \cdot \frac{\sin \angle BMP}{\sin \angle CMP}\] (by the Law of Sines in $\triangle MZY$)
 \[= \frac{\sin \angle PZM}{\sin \angle PYM} \cdot \frac{\sin \angle BMP}{\sin \angle CMP}\]
 \[= \frac{\sin \angle PCM}{\sin \angle PBM} \cdot \frac{\sin \angle BMP}{\sin \angle CMP}\]
 \[= \frac{MB}{MC} \cdot \frac{\sin \angle BMP}{\sin \angle CMP}\] (by the Law of Sines in $\triangle MBC$)
 \[= \frac{PB}{PC}\] by the Ratio Lemma.The proof is complete.

 



\textbf{Solution 3}

Use directed angles modulo $\pi$.

 Lemma. $\angle{XRY} \equiv \angle{XQZ}.$

 Proof. 
 \[\angle{XRY} \equiv \angle{XRA} - \angle{YRA} \equiv \angle{XQA} + \angle{YRB} \equiv \angle{XQA} + \angle{CPY} = \angle{XQA} + \angle{AQZ} = \angle{XQZ}.\]

 Now, it follows that (now not using directed angles)
 \[\frac{XY}{YZ} = \frac{\frac{XY}{\sin \angle{XRY}}}{\frac{YZ}{\sin \angle{XQZ}}} = \frac{\frac{RY}{\sin \angle{RXY}}}{\frac{QZ}{\sin \angle{QXZ}}} = \frac{BP}{PC}\]using the facts that $ARY$ and $APB$, $AQZ$ and $APC$ are similar triangles, and that $\frac{RA}{\sin \angle{RXA}} = \frac{QA}{\sin \angle{QXA}}$ equals twice the circumradius of the circumcircle of $AQR$.

 



\textbf{Solution 4}

We will use some construction arguments to solve the problem. Let $\angle BAC=\alpha,$ $\angle ABC=\beta,$ $\angle ACB=\gamma,$ and let $\angle APB=\theta.$ We construct lines through the points $Q,$ and $R$ that intersect with $\triangle ABC$ at the points $Q$ and $R,$ respectively, and that intersect each other at $T.$ We will construct these lines such that $\angle CQV=\angle ARV=\theta.$

 
Now we let the intersections of $AP$ with $RV$ and $QU$ be $Y'$ and $Z',$ respectively. This construction is as follows.
                    \begin{figure}[H]
                        \centering
                        \includegraphics[width=0.5\textwidth]{Solution/USAJMO_I_2013_3_Solution_Figure_2.png}
                    \end{figure}
                    

 We know that $\angle BRY'=180^\circ-\angle ARY'=180^\circ-\theta.$ Hence, we have,
 \begin{align*} \angle BRY'+\angle BPY' &=180^\circ-\theta+\theta\\ &=180^\circ. \end{align*}


 Since the opposite angles of quadrilateral $RY'PB$ add up to $180^\circ,$ it must be cyclic. Similarly, we can also show that quadrilaterals $CQZ'P,$ and $AQTR$ are also cyclic.

 Since points $Y'$ and $Z'$ lie on $AP,$ we know that,
 \[Y'=\omega_B\cap AP\]and that
 \[Z'=\omega_C\cap AP.\]

 Hence, the points $Y'$ and $Z'$ coincide with the given points $Y$ and $Z,$ respectively.

 Since quadrilateral $AQTR$ is also cyclic, we have,
 \begin{align*} \angle Y'TZ' &=180^\circ-\angle RTQ\\ &=180^\circ-(180^\circ-\angle RAQ)\\ &=\angle RAQ\\ &=\alpha. \end{align*}


 Similarly, since quadrilaterals $CQZ'P,$ and $AQTR$ are also cyclic, we have,
 \begin{align*} \angle TY'Z' &=180^\circ-\angle RY'P\\ &=180^\circ-(180^\circ-\angle RBP)\\ &=\angle RBP\\ &=\beta, \end{align*}
and,
 \begin{align*} \angle Y'Z'T &=180^\circ-\angle PZ'Q\\ &=180^\circ-(180^\circ-\angle PCQ)\\ &=\angle PCQ\\ &=\gamma. \end{align*}


 Since these three angles are of $\triangle TY'Z',$ and they are equal to corresponding angles of $\triangle ABC,$ by AA similarity, we know that $\triangle TY'Z'\sim \triangle ABC.$

 We now consider the point $X=\omega_c\cap AC.$ We know that the points $A,$ $Q,$ $T,$ and $R$ are concyclic. Hence, the points $A,$ $T,$ $X,$ and $R$ must also be concyclic.

 Hence, quadrilateral $AQTX$ is cyclic.

 
                    \begin{figure}[H]
                        \centering
                        \includegraphics[width=0.5\textwidth]{Solution/USAJMO_I_2013_3_Solution_Figure_3.png}
                    \end{figure}
                    

 Since the angles $\angle ART$ and $\angle AXT$ are inscribed in the same arc $\overarc{AT},$ we have,
 \begin{align*} \angle AXT &=\angle ART\\ &=\theta. \end{align*}


 Consider by this result, we can deduce that the homothety that maps $ABC$ to $TY'Z'$ will map $P$ to $X.$ Hence, we have that,
 \[Y'X/XZ'=BP/PC.\]

 Since $Y'=Y$ and $Z'=Z$ hence,
 \[YX/XZ=BP/PC,\]

 as required.

 



\textbf{Solution 5 (Simple Rotational Homotethy)}


                    \begin{figure}[H]
                        \centering
                        \includegraphics[width=0.5\textwidth]{Solution/USAJMO_I_2013_3_Solution_Figure_4.png}
                    \end{figure}
                    

 We begin again by noting that the three circumcircles intersect at point $M$ by Miquel's theorem. In addition, we state that the angle $\angle MQC = \alpha$, hence $\angle MPC = \angle MZC = 180 - \alpha$, as well as $\angle AQM$, from which follows that $\angle ARM = \alpha$, so $\angle BRM = 180 - \alpha$, and $\angle BPM = \alpha$. We shall prove that the triangles $\triangle MZC$ and $\triangle XMP$ and $\triangle YMB$ are similar, which will imply a rotational homotethy with angle $\angle APB$ about the point $M$, that takes $Y,X,Z$ to $B,P,C$, thus proving the problem. (In essence, just imagine we rotate $YMZX$ around M and lengthen things out and get $MCPB$ - the ratios will remain identical.)

 We do this by angle chasing. Denote angle $\angle ZMP = \beta$. From the angles labeled before, we now know $\angle ZPB = \alpha-\beta = \angle ZMC$. In addition, $\angle ZPM = \angle ZCM$. So the angles in triangle $\triangle ZMC$  are $\alpha-\beta$, $\beta$ and thus $180-\alpha$, with $\alpha-\beta$ at the point $M$. In addition, $\angle YRM = \angle YPM = \beta$, so $\angle ARY = \alpha-\beta$, so $\angle YMB = \alpha - \beta$. Since $\angle YBM = \angle YPM = \beta$, the triangle $\triangle YMB$ has angles $\alpha-\beta, \beta, 180-\alpha$ also, with $\alpha-\beta$ at $M$. Finally, In triangle $\triangle XPM$, we already know the angle $\angle XPM$ to be $\beta$; we also can find that $\angle AQM = 180-\alpha$, so $\angle AXM = \alpha$, so $\angle MXP = 180-\alpha$, so $\angle XMP = \alpha-\beta$. Thus, the three triangles are similar have a common point $M$, which proves that there is a rotational homotethy around $M$ that maps $Y, X, Z$ to $B, P, C$ as desired. 

 \textit{Solution by SimilarTriangle.}

 


	\item (\textit{USAJMO 2013}) Let $f(n)$ be the number of ways to write $n$ as a sum of powers of $2$, where we keep track of the order of the summation.  For example, $f(4)=6$ because $4$ can be written as $4$, $2+2$, $2+1+1$, $1+2+1$, $1+1+2$, and $1+1+1+1$.  Find the smallest $n$ greater than $2013$ for which $f(n)$ is odd.


 



\textbf{Solution}

First of all, note that $f(n)$ = $\sum_{i=0}^{k} f(n-2^{i})$ where $k$ is the largest integer such that $2^k \le n$.  We let $f(0) = 1$ for convenience.  

 From here, we proceed by induction, with our claim being that the only $n$ such that $f(n)$ is odd are $n$ representable of the form $2^{a} - 1, a \in \mathbb{Z}$

 We induct on $a$.  It is trivially true for $a = 0$ and $a = 1$.  From here, we show that, if the only numbers $n \le 2^{a-1} - 1$ where $f(n)$ is odd are of the form described above, then the only numbers $n \le 2^{a} -1$ that are odd are of that form.  We first consider all numbers $b$, such that $2^{a-1} \le b \le 2^{a} - 2$, going from the lower bound to the upper bound (a mini induction, you might say).  We know that $f(b) = \sum_{i=0}^{a-1} f(b-2^{i})$.  For a number in this summation to be odd, $b - 2^i = 2^m -1 \rightarrow b = 2^i + 2^m - 1$.  However, we know that $b > 2^{a-1}$, so $m$ must be equal to $a-1$, or else $b$ cannot be in that interval.  Now, from this, we know that $i < a-1$, as $b<2^{a} - 1$.  Therefore, $i$ and $m$ are distinct, and thus $f(b - 2^i)$ and $f(b- 2^{a-1})$ are odd; since there are just two odd numbers, the ending sum for any $b$ is even.  Finally, considering $2^{a} - 1$, the only odd number is $f(2^{a} - 1 - 2^{a-1})$, so the ending sum is odd.  $\Box$

 The smallest $n$ greater than $2013$ expressible as $2^d - 1, d \in \mathbb{N}$ is $2^{11} -1 = \boxed{2047}$

 



\textbf{Solution with Explanations}

Of course, as with any number theory problem, use actual numbers to start, not variables! By plotting out the first few sums (do it!) and looking for patterns, we observe that $f(n)=\sum_{\textrm{power}=0}^{\textrm{pow}_{\textrm{larg}}} f(n-2^{\textrm{power}})$, where $\textrm{pow}_{\textrm{larg}}$ represents the largest power of $2$ that is smaller than $n$. I will call this sum the Divine Sign, or DS.

 But wait a minute... we are trying to determine odd/even of $f(n)$. Why not call all the evens 0 and odds 1, basically using mod 2? Sounds so simple. Draw a small table for the values: as $n$ goes up from $0$, you get: $1,1,0,1,0,0,0,1,0...$. We have to set $f(0)=1$ for this to work. Already it looks like $f(n)$ is only odd if $n=2^{\textrm{power}}-1$.

 The only tool here is induction. The base case is clearly established. Then let's assume we successfully made our claim up to $2^n-1$. We need to visit numbers from $2^n$ to $2^{n+1}-1$. Realize that $2^n$ has $0$ for $f$ because there will be two numbers in DS that give a $f$ of one: $2^{n-1}$ and $1$. 

 But to look at whether a value of $f(\textrm{number})$ is 1 or 0, we need to revisit our first equation. We can answer this rather natural question: When will a number to be inducted upon, say $2^n+k$, ever have a 1 as $f(\textrm{number})$ in the DS equation? Well- because by our assumption of the claim up to $2^n-1$, we know that the only way for that to happen is if $2^n+k-2^{\textrm{power}}$ in the DS is equal to $2^{\textrm{Some power}} - 1$. Clearly $1 \leq k \leq 2^n - 1$.

 Finally, we can simplify. Using our last equation, $2^n+k-2^{\textrm{power}}=2^{\textrm{Some power}}-1$, regrouping gives $2^n+k=2^{\textrm{power}}+2^{\textrm{Some power}}-1$. 

 Most importantly, realize that $\textrm{power}$ can be from $0$ to $n$, because of the restraints on $k$ mentioned earlier. Same with $\textrm{Some power}$. Immediately at least one of $\textrm{power}$ and $\textrm{Some power}$ has to be $n$. If both were smaller, LHS is greater, contradiction. If both were greater, RHS is greater, contradiction.

 Therefore, by setting one of $\textrm{power}$ or $\textrm{Some power}$ to $n$, we realize $k=2^{\textrm{A certain power}}-1$.

 The conclusion is clear, right? Each $k$ from $1$ to $2^{n-1}-1$ yields two distinct cases: one of $\textrm{power}$ and $\textrm{Some power}$ is equal to $n$, while the other is LESS THAN $n$. But for $k=2^n-1$, there is ONE CASE: BOTH values have to equal $n$. Therefore, the only $k$ that has $f(2^n+k)$ as odd must only be $2^n-1$, because the other ones yield a $f$ of 1+1=0 in our mod. That proves our induction for a new power of 2, namely $n+1$, meaning that $f(\textrm{number})$ is only odd if $\textrm{number} = 2^{\textrm{Power of two}} - 1$, and we are almost done...

 Thus, the answer is $2^{11}-1=\boxed{2047}$.

 This was pretty intuitive, and realizing quite $elementary$ steps about how powers of 2 work gave us the solution! Estimated time: ~50 minutes. 

 minor edits ~mathboy100

 


	\item (\textit{USAJMO 2013}) Quadrilateral $XABY$ is inscribed in the semicircle $\omega$ with diameter $XY$.  Segments $AY$ and $BX$ meet at $P$.  Point $Z$ is the foot of the perpendicular from $P$ to line $XY$.  Point $C$ lies on $\omega$ such that line $XC$ is perpendicular to line $AZ$.  Let $Q$ be the intersection of segments $AY$ and $XC$.  Prove that 
 \[\dfrac{BY}{XP}+\dfrac{CY}{XQ}=\dfrac{AY}{AX}.\]


 



\textbf{Solution}

Using the Law of Sines and simplifying, we have 
 \[\frac{BY}{XP}+\frac{CY}{XQ}=\frac{\sin \angle PXY \sin \angle XPY+\sin \angle QXY\sin \angle XQY}{\sin \angle AYX}.\]

 It is easy to see that $APZX$ is cyclic. Also, we are given $XQ\perp AZ$. Then we have 
 \begin{align*} \frac{\sin \angle PXY\sin \angle XPY+\sin \angle QXY\sin \angle XQY}{\sin \angle AYX} &= \frac{\sin \angle YAZ\sin \angle XZA+\cos \angle XZA\cos \angle YAZ}{\sin \angle AYX} \\ &= \frac{\cos(\angle XZA-\angle YAZ)}{\sin \angle AYX} \\ &= \frac{\cos \angle AYX}{\sin \angle AYX} \\ &= \cot \angle AYX \\ &= \frac{AY}{AX}, \end{align*}
 and we are done.

 



\textbf{Solution 1}

Let us use coordinates. Let O, the center of the circle, be (0,0). WLOG the radius of the circle is 1, so set Y (1,0) and X (-1,0). Also, for arbitrary constants $a$ and $b$ set A $(\cos a, \sin a)$ and B $(\cos b, \sin b)$. Now, let's use our coordinate tools. It is easily derived that the equation of $BX$ is $y = \frac{\sin b}{1 + \cos b}(x + 1) = v(x+1)$ and the equation of $AY$ is $y = \frac{\sin a}{1 - \cos a}(x - 1) = u(x-1)$, where $u$ and $v$ are defined appropriately. Thus, by equating the y's in the equation we find the intersection of these lines, $P$, is $\left(\frac{u-v}{u+v}\right), \frac{2uv}{u+v})$. Also, $Z\left(\frac{u-v}{u+v}\right), 0)$. It shall be left to the reader to find the slope of $AZ$, the coordinates of Q and C, and use the distance formula to verify that $\frac{BY}{XP} + \frac{CY}{XQ} = \frac{AY}{AX}$.

 



\textbf{Solution 2}

First of all 

 
 \[\angle BXY = \angle PAZ =\angle AXQ =\angle AXC\] since the quadrilateral $APZX$ is cyclic, and triangle $AXQ$ is rectangle, and $CX$ is orthogonal to $AZ$. Now 

 
 \[\angle BXY =\angle BAY =\angle AXC\] because $XABY$ is cyclic and we have proved that 

 
 \[\angle AXC = \angle BXY\] so $BC$ is parallel to $AY$, and 
 \[AC=BY, CY=AB\] Now by Ptolomey's theorem on $APZX$ we have 
 \[(AX)(PZ)+(AP)(XZ)=(AZ)(PX)\] we see that triangles $PXZ$ and $QXA$ are similar since 
 \[\angle QAX= \angle PZX= 90\] and 
 \[\angle AXC = \angle BXY\] is already proven, so 
 \[(AX)(PZ)=(AQ)(XZ)\] Substituting yields 
 \[(AQ)(XZ)+(AP)(XZ)=(AZ)(PX)\] dividing by $(PX)(XZ)$ We get 
 \[\frac {AQ+AP}{XP} = \frac {AZ}{XZ}\] Now triangles $AYZ$, and $XYP$ are similar so 
 \[\frac {AY}{AZ}= \frac {XY}{XP}\] but also triangles $XPY$ and $XZB$ are similar and we get 
 \[\frac {XY}{XP}= \frac {XB}{XZ}\] Comparing we have, 
 \[\frac {AY}{XB}= \frac {AZ}{XZ}\] Substituting, 
 \[\frac {AQ+AP}{XP}= \frac {AY}{XB}\] Dividing the new relation by $AX$ and multiplying by $XB$ we get 
 \[\frac{XB(AQ+AP)}{(XP)(AX)} = \frac {AY}{AX}\] but 
 \[\frac {XB}{AX}= \frac {XY}{XQ}\] since triangles $AXB$ and $QXY$ are similar, because 
 \[\angle AYX= \angle ABX\] and 
 \[\angle AXB= \angle CXY\] since $CY=AB$ Substituting again we get 
 \[\frac {XY(AQ)+XY(AP)}{(XP)(XQ)} =\frac {AY}{AX}\] Now since triangles $ACQ$ and $XYQ$ are similar we have 
 \[XY(AQ)=AC(XQ)\] and by the similarity of $APB$ and $XPY$, we get 
 \[AB(CP)=XY(AP)\] so substituting, and separating terms we get 
 \[\frac{AC}{XP} + \frac{AB}{XQ} = \frac{AY}{AX}\] In the beginning we prove that $AC=BY$ and $AB=CY$ so 
 \[\frac{BY}{XP} + \frac{CY}{XQ} = \frac{AY}{AX}\] $\blacksquare$

 



\textbf{Solution 3}

It is obvious that 
 \[\angle AXB=\angle CXY=\alpha\]for some value $\alpha$. Also, note that $\angle BYA=\alpha$. Set 
 \[\angle BXC=\angle BYC=\beta.\] We have
 \[\frac{XC}{CY}=\tan {\angle CYZ}=\tan (90-\alpha)\]and 
 \[\frac{CQ}{CY}=\tan {\angle CYQ}=\tan (\alpha+\beta).\]This gives
 \[\frac{CY}{XQ}=\frac{1}{\tan (90-\alpha)-\tan (\alpha+\beta)}.\]Similarly, we can deduce that 
 \[\frac{BY}{XP}=\frac{1}{\tan (90-\alpha-\beta)-\tan {\alpha}}.\]Adding gives
 \[\frac{\tan \alpha}{1-\tan (\alpha+\beta)\tan {\beta}}+\frac{\tan (\alpha+\beta)}{1-\tan (\alpha+\beta)\tan {\beta}}=\tan (2\alpha +\beta)=\frac{AY}{AX}.\]

 


 



\textbf{Solution 4}

First, since $XY$ is the diameter and $A$, $B$, and $C$ lie on the circle, 
 \[\angle {XAY} = \angle {XBY} = \angle{XCY} = 90\]. Next, because $AZ$ and $CY$ are both perpendicular to $CX$, we have $AZ$ to be parallel to $CY$. 

 Now looking at quadrilateral $APZX$, we see that this is cyclic because 
 \[\angle {PAX} + \angle {PZX} = 90+90 = 180.\] Set $\alpha = \angle{BXA} = \angle{BYA}$, and $\beta = \angle{BXC} = \angle{BYC}$. Now, 
 \[\angle{AYC} = \angle{YAZ}\] since $AZ$ and $CY$ are parallel. Also, 
 \[\angle{PAZ} = \angle{PXZ} = \alpha + \beta.\] That means 
 \[\angle{PXZ} = \angle{PXQ} + \angle{QXZ} = \beta + \angle{QXZ} = \alpha + \beta\] so 
 \[\angle{QXZ} = \alpha.\] This means $\angle{QXZ} = \angle{YBC} = \alpha$, so $BC$ and $AY$ are parallel. Finally, we can look at the equation. We know 
 \[XP\cos{\alpha} = AX,\] so $XP = \frac{AX}{\cos{\alpha}}.$We also know 
 \[XQ\cos(\alpha+\beta) = AX,\] so $XQ = \frac{AX}{\cos(\alpha+\beta)}.$Plugging this into the LHS of the equation, we get 
 \[\frac{BY\cos\alpha}{AX}+\frac{CY\cos(\alpha+\beta)}{AZ}.\]Now, let $H$ be the point on $AY$ such that $BH$ is perpendicular to $AY$. Also, since $\angle{AYB} = \angle{CXY}$, their arcs have equal length, and $AB=CY$. Now, the LHS is simplified even more to 
 \[\frac{AB\cos\alpha}{AX}+\frac{CY\cos(\alpha+\beta)}{AZ}\] which is equal to 
 \[\frac{AH+YH}{AZ}\] which is equal to 
 \[\frac{AY}{AX}.\] This completes the proof. 

 


 


	\item (\textit{USAJMO 2013}) Find all real numbers $x,y,z\geq 1$ satisfying 
 \[\min(\sqrt{x+xyz},\sqrt{y+xyz},\sqrt{z+xyz})=\sqrt{x-1}+\sqrt{y-1}+\sqrt{z-1}.\]


 



\textbf{Solution with Thought Process}

Without loss of generality, let $1 \le x \le y \le z$. Then $\sqrt{x + xyz} = \sqrt{x - 1} + \sqrt{y - 1} + \sqrt{z - 1}$.

 Suppose x = y = z. Then $\sqrt{x + x^3} = 3\sqrt{x-1}$, so $x + x^3 = 9x - 9$. It is easily verified that $x^3 - 8x + 9 = 0$ has no solution in positive numbers greater than 1. Thus, $\sqrt{x + xyz} \ge \sqrt{x - 1} + \sqrt{y - 1} + \sqrt{z - 1}$ for x = y = z. We suspect if the inequality always holds.

 Let x = 1. Then we have $\sqrt{1 + yz} \ge \sqrt{y-1} + \sqrt{z-1}$, which simplifies to 
 \[1 + yz \ge y + z - 2 + 2\sqrt{(y-1)(z-1)}\] and hence 
 \[yz - y - z + 3 \ge 2\sqrt{(y-1)(z-1)}\] Let us try a few examples: if y = z = 2, we have $3 > 2$; if y = z, we have $y^2 - 2y + 3 \ge 2(y-1)$, which reduces to $y^2 - 4y + 5 \ge 0$. The discriminant (16 - 20) is negative, so in fact the inequality is strict. Now notice that yz - y - z + 3 = (y-1)(z-1) + 2. Now we see we can let $u = \sqrt{(y-1)(z-1)}$! Thus, 
 \[u^2 - 2u + 2 = (u-1)^2 + 1 > 0\] and the claim holds for x = 1.

 If x > 1, we see the $\sqrt{x - 1}$ will provide a huge obstacle when squaring. But, using the identity $(x+y+z)^2 = x^2 + y^2 + z^2 + xy + yz + xz$:
 \[x + xyz \ge x - 1 + y - 1 + z - 1 + 2\sqrt{(x-1)(y-1)} + 2\sqrt{(y-1)(z-1)} + 2\sqrt{(x-1)(y-1)}\]which leads to
 \[xyz \ge y + z - 3 + 2\sqrt{(x-1)(y-1)} + 2\sqrt{(y-1)(z-1)} + 2\sqrt{(x-1)(z-1)}\]Again, we experiment. If x = 2, y = 3, and z = 3, then $18 > 7 + 4\sqrt{6}$.

 Now, we see the finish: setting $u = \sqrt{x-1}$ gives $x = u^2 + 1$. We can solve a quadratic in u! Because this problem is a \#6, the crown jewel of USAJMO problems, we do not hesitate in computing the messy computations:
 \[u^2(yz) - u(2\sqrt{y-1} + 2\sqrt{z-1}) + (3 + yz - y - z - 2\sqrt{(y-1)(z-1)}) \ge 0\]

 Because the coefficient of $u^2$ is positive, all we need to do is to verify that the discriminant is nonpositive:
 \[b^2 - 4ac = 4(y-1) + 4(z-1) - 8\sqrt{(y-1)(z-1)} - yz(12 + 4yz - 4y - 4z - 8\sqrt{(y-1)(z-1)})\]

 Let us try a few examples. If y = z, then the discriminant D = $8(y-1) - 8(y-1) - yz(12 + 4y^2 - 8y - 8(y-1)) = -yz(4y^2 - 16y + 20) = -4yz(y^2 - 4y + 5) < 0$.

 We are almost done, but we need to find the correct argument. (How frustrating!)Success! The discriminant is negative. Thus, we can replace our claim with a strict one, and there are no real solutions to the original equation in the hypothesis.

 --Thinking Process by suli

 


\end{enumerate}